\documentclass[11pt]{article}
\usepackage[margin=1in]{geometry}
\usepackage{amsmath,amssymb,amsthm}
\usepackage{hyperref}
\hypersetup{hidelinks}

\newtheorem{theorem}{Theorem}
\newtheorem{lemma}[theorem]{Lemma}
\newtheorem{corollary}[theorem]{Corollary}
\theoremstyle{definition}
\newtheorem{definition}{Definition}

\setlength{\parskip}{4pt plus 1pt minus 1pt}
\setlength{\parindent}{0pt}

\title{Evaluation-Driven State Revision and the\\Encoding Hierarchy Under Bounded Context}
\author{Patrick D.\ McCarthy}
\date{}

\begin{document}
\maketitle

\begin{abstract}
We prove two linked results for any entity persisting under bounded
active context with survival pressure. First, agency ($\gamma > 0$)
necessarily produces evaluation-driven state revision---monotone
descent on expected uncertainty $H(W \mid K)$---with convergence
guaranteed by a scalar Lyapunov argument: $H(W \mid K) \geq 0$ is
bounded below, each informative update reduces $H(W \mid K)$ by
exactly the information gain $G(a, K)$, and monotone convergence
applies without dimensional decomposition. The process is not gradient
descent (there is no differentiable loss surface over $K$) but shares
the descent property: each update produces lower
expected uncertainty. Descent never terminates under positive novelty
rate: as long as the entity encounters states where $U > 0$, actions
with positive information gain exist. Second, the certainty threshold for encoding an action
at level $L_i$ is $\theta_i = 1 - \varepsilon / C_i$, where $C_i$ is
the cost of failure at that level. This derives the learning rate
hierarchy: evolutionary rates are necessarily orders of magnitude lower
than individual learning rates, not by assumption but from
cost-of-failure differences.
\end{abstract}

% ------------------------------------------------------------------
\section{Introduction}
\label{sec:intro}
% ------------------------------------------------------------------

An entity with agency $\gamma > 0$ acts on the world and receives feedback.
What is the necessary structure of learning? Papers~1--3 in this
series established that any entity persisting under bounded active
context must maintain a directed graph over propositions
\cite{McCarthy2026a}, that control architectures converge to escalation at scale
\cite{McCarthy2026b}, and that knowledge and action are inseparable
under survival pressure \cite{McCarthy2026c}. What remains open is two
questions: what \emph{dynamics} govern the updating process, and what
determines the \emph{threshold} at which an action earns promotion to
a more permanent encoding level.

This paper answers both. We prove that the updating process is
monotone descent on expected uncertainty $H(W \mid K)$, with
convergence guaranteed by a scalar Lyapunov argument that requires no
dimensional decomposition (Theorem~\ref{thm:gradient}). We then prove that the certainty
threshold for encoding at level $L_i$ is determined by the cost of
failure at that level, which derives the learning rate hierarchy across
encoding levels (Theorem~\ref{thm:encoding}).

Several lines of prior work converge on these questions from different
angles. Dupoux, LeCun, and Malik \cite{DupouxLeCunMalik2026} propose
a bilevel optimization framework (evolution as outer loop, development
as inner loop) but treat the learning rate relationship between levels
as a design parameter rather than deriving it. Cover and Thomas
\cite{Cover2006} provide the information-theoretic foundations that
anchor our convergence argument. Schmidhuber \cite{Schmidhuber2010} identifies
curiosity as compression progress, structurally identical to our
learning rate $\eta_M$. Hutter \cite{Hutter2005} constructs AIXI as
an idealized agent but assumes unbounded computation, precisely the
constraint our framework takes as central. Chollet
\cite{Chollet2019} defines intelligence as skill-acquisition
efficiency, close to $\eta_M$ but treating prior knowledge as
confound rather than constitutive element. Hinton and Nowlan
\cite{Hinton1987} demonstrate that learning can guide evolution (the
Baldwin effect) without formalizing the threshold that governs
promotion. Finn et al.\ \cite{FinnEtAl2017} and Thrun
\cite{Thrun1998} develop meta-learning algorithms that operate as $M$
on $M$, the $M^{(2)}$ level in our hierarchy, but do not derive
why the meta-level necessarily learns more slowly.

% ------------------------------------------------------------------
\section{Formal Framework}
\label{sec:framework}
% ------------------------------------------------------------------

We introduce the minimal definitions required for the two main
results, extending the framework of \cite{McCarthy2026a}.

\begin{definition}[Entity, World, Uncertainty]
\label{def:entity}
An entity $E = (K, A, \sigma, \gamma)$ consists of a knowledge base $K$
(propositions about world $W$), an action space $A$ (executable
actions), a selection mechanism $\sigma$ that chooses actions given
context, and an agency parameter $\gamma$. The world state space $W$ is a
measurable space; the true state $w \in W$ is not directly observable.
The \emph{uncertainty} of $E$ about state $w$ given knowledge $K$ is:
\[
U(w, K) = -\log P(w \mid K)
\]
The expected uncertainty is the conditional entropy:
\[
H(W \mid K) = \mathbb{E}_w[U(w, K)] = -\sum_w P(w \mid K) \log P(w \mid K)
\]
\end{definition}

\begin{definition}[Agency]
\label{def:agency}
The agency parameter $\gamma \in [0,1]$ is the probability that $E$
engages its learning mechanism $M$ in response to a learning signal.
At $\gamma = 0$, $A$ and $K$ are fixed forever: the entity receives
signals but never updates. The operational signature of $\gamma > 0$ is
that the entity bounds uncertainty beyond what immediate survival
requires, freeing active context capacity for novel states. An entity
that updates only under immediate threat has $\gamma \approx 0$; an
entity that builds reserve capacity for states it has not yet
encountered has $\gamma > 0$.
\end{definition}

\begin{definition}[Action Feedback]
\label{def:feedback}
For action $a$ executed in world state $w$, the \emph{feedback}
$\varphi(a, w)$ is the observable consequence. The entity updates its
knowledge base: $K_a = K \cup \{p_a\}$, where $p_a$ is the
proposition encoding the feedback. The \emph{prediction-error signal} is:
\[
\delta = \varphi(a, w) - \mathbb{E}[\varphi(a, \cdot) \mid K]
\]
\end{definition}

\begin{definition}[Information Gain]
\label{def:gain}
When entity $E$ executes action $a$ in world state $w$, it receives
feedback $\varphi(a, w)$ encoded as proposition $p_a$
(Definition~\ref{def:feedback}). The \emph{information gain} is the
mutual information between the feedback proposition and the world
state, conditioned on current knowledge:
\[
G(a, K) = I(p_a; W \mid K) \geq 0
\]
Non-negativity follows from the definition of mutual information
\cite{Cover2006}. $G = 0$ if and only if $p_a$ provides no
information about $W$ beyond what $K$ already contains: the feedback
is redundant given current knowledge.
\end{definition}

\begin{definition}[Novelty Rate]
\label{def:novelty}
The \emph{novelty rate} $\nu(E)$ is the rate at which entity $E$
encounters world states $w$ where $U(w, K) > 0$. $\nu$ is a property
of the entity-environment coupling, not of $W$ alone: the same
environment presents different novelty rates to entities with different
$K$. $\nu(E) > 0$ whenever the entity's environment is not fully
predicted by $K$.
\end{definition}

\begin{definition}[Encoding Hierarchy]
\label{def:hierarchy}
The encoding space $\mathcal{E}(c, r)$ is a continuous space
parameterized by cost $c$ and reversibility $r$. Within this space,
interfaces $L_0, L_1, \ldots, L_n$ define discrete encoding levels
with increasing cost and reversibility:
\begin{itemize}
\item $L_0$ (genomic): highest cost of error, lowest reversibility,
   slowest update rate.
\item $L_n$ (active context): lowest cost of error, highest
   reversibility, fastest update rate.
\end{itemize}
Escalation \cite{McCarthy2026b} pushes proven actions
upward toward more permanent levels. The present paper derives the
threshold that governs this promotion.
\end{definition}

% ------------------------------------------------------------------
\section{Monotone Descent on Uncertainty}
\label{sec:descent}
% ------------------------------------------------------------------

\begin{theorem}[Evaluation-Driven State Revision]
\label{thm:gradient}
Let $E = (K, A, \sigma, \gamma)$ be an entity with $\gamma > 0$ persisting in
world $W$ under bounded active context $C_n$.

\textbf{Part I.} Every executed action generates a prediction-error
signal, and $\gamma > 0$ guarantees engagement with that signal,
producing evaluation-driven state revision---monotone descent on
expected uncertainty $H(W \mid K)$---with convergence to a local
minimum.

\textbf{Part II.} Under positive novelty rate $\nu(E) > 0$
(Definition~\ref{def:novelty}), descent never terminates globally:
the entity continues to encounter states where $U(w,K) > 0$, so
descent never terminates.
\end{theorem}

\begin{proof}
\textbf{Part I: Forced Signal, Response Gate, and Convergence.}

Every executed action $a$ generates feedback $\varphi(a, w)$ from the
world. This signal arrives whether or not the entity has $\gamma > 0$: the
world responds to actions regardless of the actor's capacity to learn
from them. The prediction error $\delta = \varphi(a, w) -
\mathbb{E}[\varphi(a, \cdot) \mid K]$ is therefore produced at every
action.

Agency $\gamma$ is the response gate. An entity with $\gamma = 0$ receives
every signal $\delta$ but never updates: $(A, K)$ remain static. An
entity with $\gamma > 0$ engages $M$ with probability $\gamma$ on each
signal, integrating the feedback into $K$ by forming
$K' = K \cup \{p_a\}$, where $p_a$ is the proposition encoding the
feedback. $K$ is a set of propositions, not a parameter vector, so
$\nabla_K U$ is undefined: there is no differentiable loss surface over
$K$. The descent property is instead established directly via the
Lyapunov function $H(W \mid K) = \mathbb{E}_w[U(w, K)]$, the expected
uncertainty of $K$ about $W$.

We now establish convergence via a scalar Lyapunov argument that
requires no dimensional decomposition.

\emph{Step 1: Each informative update strictly reduces $H(W \mid K)$.}
When the entity encounters a novel state and executes action $a$, the
feedback $p_a$ is informative: $G(a, K) = I(p_a; W \mid K) > 0$
because the feedback resolves uncertainty that $K$ did not already
contain \cite{Cover2006}. When $M$ integrates this feedback, forming
$K' = K \cup \{p_a\}$:
\[
H(W \mid K') = H(W \mid K) - G(a, K)
\]
This identity follows from the chain rule for mutual information under
the assumption that the entity performs exact Bayesian conditioning:
$H(W \mid K') = H(W \mid K, p_a)$, i.e., adding $p_a$ to $K$ is
equivalent to conditioning on $p_a$. This is an idealization---any
bounded entity is an approximate updater---but the approximation is
conservative: an approximate updater achieves $H(W \mid K') \leq
H(W \mid K)$ with $G(a,K)$ as an upper bound on the reduction, which
suffices for the monotone convergence argument below. The exact
Bayesian case is the cleanest statement; the result degrades gracefully
under approximation.
Since $G(a, K) > 0$ for novel encounters, $H(W \mid K') < H(W \mid K)$
strictly. The Lyapunov function decreases at every informative update.

\emph{Step 2: Monotone convergence.} Consider the sequence
$\{H_t\} = \{H(W \mid K(t))\}$ indexed by informative update steps
$t$ (non-novel encounters produce no update and leave $H_t$
unchanged). By Step~1, $\{H_t\}$ is strictly decreasing at each
informative step. By definition, $H(W \mid K) \geq 0$, so the
sequence is bounded below. By the monotone sequence theorem (a bounded monotone sequence of reals
converges),
$\{H_t\}$ converges. Since the sequence is non-increasing and
convergent:
\[
\sum_{t=0}^{\infty} G(a_t, K(t)) \leq H(W \mid K(0)) < \infty
\]
Therefore $G(a_t, K(t)) \to 0$ as $t \to \infty$. The entity
converges to a state where available actions yield diminishing
informational returns.

\emph{Step 3: Dimensionality irrelevance.} The entire argument
operates on the scalar $H(W \mid K)$. The entity's knowledge base may
grow as new propositions are added. Each new proposition $p_a$ from an
informative encounter reduces $H(W \mid K)$ by exactly $G(a, K)$
(Step~1). Cross-grounding between propositions accelerates convergence
by allowing a single observation to update multiple related beliefs.
The Lyapunov bound $H(W \mid K) \geq 0$ holds regardless of the
size of $K$ or the parameterization of $M$.

\textbf{Part II: Non-Termination Under Positive Novelty.}

When $\nu(E) > 0$ (Definition~\ref{def:novelty}), the entity
encounters states $w$ with $U(w, K(t)) > 0$ at positive rate. Each
such encounter produces an action $a$ with $G(a, K(t)) > 0$: the
feedback from acting under uncertainty is informative. Therefore at
every finite $t$ there exists some $a'$ with $G(a', K(t)) > 0$.

Descent never terminates globally. Descent converges locally
(Part~I) but never terminates globally: as long as the environment
presents novel states, there remain actions with positive information
gain.

Non-termination does not require unbounded world complexity. A finite
world with $\nu(E) > 0$ suffices: even if $W$ is finite, the entity
may never fully predict it because new evidence can invalidate or
refine existing propositions, and the combinatorial space of
entity-environment interactions exceeds what $K$ can index under
bounded $C_n$.

\emph{Interaction with the curation regime.} The convergence argument
in Part~I assumes $K$ grows monotonically ($K' = K \cup \{p_a\}$). In
the curation regime established in \cite{McCarthy2026c}, $|K|$ is
bounded at $K_{\max}$ and every addition requires displacement. Does
displacement break convergence?

The Lyapunov function must be refined. Raw $H(W \mid K)$ is the wrong
quantity in the curation regime: displacing a proposition about an
irrelevant domain in favor of one about a survival-critical domain
could increase $H(W \mid K)$ overall while improving the entity's
survival prospects. The correct Lyapunov function is
\emph{survival-weighted expected uncertainty}:
\[
H_{\mathcal{T}}(W \mid K) = \mathbb{E}_{w \sim \mathcal{T}}[U(w, K)]
\]
where $\mathcal{T}$ is the entity's task distribution over
survival-relevant world states. This is the expected uncertainty over
the states that actually matter for persistence.

Under the cost-benefit displacement criterion of the $K/A$
inseparability theorem \cite{McCarthy2026c}, proposition $p$ is
displaced by $q$ only when
\[
  \mathbb{E}[B(q, \mathcal{T})] - c_{\text{retain}}(q, K)
  \;>\;
  \mathbb{E}[B(p, \mathcal{T})] - c_{\text{retain}}(p, K),
\]
where $B$ is expected survival benefit. Since $B$ is defined over
$\mathcal{T}$, this criterion directly reduces $H_{\mathcal{T}}$: the
entity replaces coverage of less survival-relevant states with coverage
of more survival-relevant states. $H_{\mathcal{T}}(W \mid K)$ is non-increasing across displacements.

$\mathcal{T}$ itself is non-stationary: as the entity learns, moves,
and the environment shifts, the distribution of survival-relevant
states changes. Under non-stationary $\mathcal{T}$, the Lyapunov
convergence argument of Part~I does not apply in its global form: a
Lyapunov argument requires a fixed function that decreases along
trajectories, and $H_{\mathcal{T}_t}$ at step $t$ and
$H_{\mathcal{T}_{t+1}}$ at step $t+1$ evaluate different functions
when $\mathcal{T}$ drifts. What the curation regime establishes is
\emph{local, instantaneous} descent: each displacement reduces
$H_{\mathcal{T}}$ evaluated against the current $\mathcal{T}$ at the
moment of displacement. This is a weaker claim than global
convergence, but it is the operationally relevant one. The entity
cannot plan against future task distributions it has not yet
encountered; it can only improve its coverage of the tasks it currently
faces. The non-stationarity of $\mathcal{T}$ is precisely what drives
non-termination (Part~II): the target keeps moving, so the entity
never reaches $H_{\mathcal{T}} = 0$ globally, even as it descends
locally at every step.

\textbf{Remark.} An entity with $\gamma = 0$ never updates. As $W$
changes, its fixed $K$ becomes increasingly inaccurate: $U(w_t, K)$
grows until it exceeds the lethality threshold $U_{\mathrm{lethal}}$.
Selection removes such entities. Therefore $\gamma > 0$ is necessary for
persistence, a claim established fully in Paper~5.
\end{proof}

% ------------------------------------------------------------------
\section{Multi-Timescale Structure}
\label{sec:multiscale}
% ------------------------------------------------------------------

\begin{corollary}[Individual Learning and Evolution Share Update Structure]
\label{cor:multiscale}
Both individual learning and evolutionary selection instantiate the
same abstract update: execute, receive signal, retain or revise. They
differ in mechanism, object, and timescale.
\end{corollary}

\begin{proof}
\textbf{Individual learning.} The entity updates its knowledge and
action space directly from continuous prediction-error signals
$\delta$. Each update revises $K$ in a direction that reduces
$\mathbb{E}[U]$ (Theorem~\ref{thm:gradient}, Part~I). The update is
directed: $\delta$ specifies both magnitude and direction of revision.
The learning rate $\eta_{\mathrm{ind}}$ is determined by the capacity
of mechanism $M$.

\textbf{Evolutionary selection.} The population-level update operates
on a different object (the heritable component of $A$) via a different
mechanism:
\[
\Delta a_{\mathrm{species}} = \mathrm{select}(\mathrm{mutate}(A_{\mathrm{population}}))
\]
Mutation is undirected: random perturbation of heritable parameters.
Selection is binary: survive or not. No individual organism descends a
gradient at the evolutionary timescale; the population distribution
shifts toward lower $U$ through differential survival.

Both processes move the entity (individual or lineage) toward lower
$U$ at their respective timescales. Evaluation-driven revision is the
efficient limit: directed $\delta$ combined with capable $M$ produces
the fastest feasible descent. Evolution is the general case: random
perturbation combined with selection produces descent without requiring
any individual to compute a prediction error.

The two are not independent. Individual learning discovers actions at
$L_n$ that, if sufficiently validated, promote to $L_0$ via the
encoding hierarchy (Theorem~\ref{thm:encoding}). This is the Baldwin
effect \cite{Hinton1987}, here derived rather than observed.
\end{proof}

% ------------------------------------------------------------------
\section{Encoding Permanence}
\label{sec:encoding}
% ------------------------------------------------------------------

\begin{theorem}[Encoding Permanence Scales with Rigor]
\label{thm:encoding}
Let $L_0, \ldots, L_n$ be encoding levels with associated failure
costs $C_0 > C_1 > \cdots > C_n$ and reversibility costs
$r_0 > r_1 > \cdots > r_n$. Let $\varepsilon$ be the maximum
acceptable expected error cost. Then the minimum certainty threshold
for encoding an action at level $L_i$ is:
\[
\theta_i = 1 - \frac{\varepsilon}{C_i}
\]
\end{theorem}

\begin{proof}
Encoding action $a$ at level $L_i$ has reversibility cost $r_i$:
correcting a bad encoding is not impossible but expensive in proportion
to $r_i$ (at $L_0$, correction requires generations of selection; at
$L_n$, a single update suffices). If the encoding is erroneous---$a$
is not in fact reliably adaptive---the cost of that error is $C_i$. If
$\theta$ is the certainty that $a$ is correct, then the expected error
cost is:
\[
\mathbb{E}[\text{error cost}] = (1 - \theta) \cdot C_i
\]
Encoding is safe when the expected error cost does not exceed the
acceptable threshold:
\[
(1 - \theta) \cdot C_i \leq \varepsilon
\]
Rearranging:
\[
\theta \geq 1 - \frac{\varepsilon}{C_i}
\]
The minimum safe threshold is therefore $\theta_i = 1 -
\varepsilon / C_i$. This is the break-even point on a cost continuum:
the certainty at which the expected benefit of encoding at $L_i$
exceeds the expected cost of error weighted by $C_i$.

\textbf{Consequence.} At $L_0$ (genomic encoding), $C_0$ is the cost
of lineage extinction, orders of magnitude larger than $\varepsilon$.
Therefore $\theta_0 \approx 1$: near-certainty is required before the
expected benefit justifies the risk. A bad encoding at $L_0$ does not
kill one individual; it kills the lineage.
At $L_n$ (active context), $C_n$ is small (cost of a single failed
action), so $\theta_n$ is meaningfully less than $1$: the expected
benefit of trying exceeds the expected cost of failure at moderate
confidence. Rigor is proportional to reversibility cost.
\end{proof}

\textbf{Error propagation through the hierarchy.} A natural objection
holds that errors compound across levels: if each level has error rate
$1 - \theta_i$, the cumulative error grows multiplicatively. This
objection is based on a naive cascading model. Under escalation
\cite{McCarthy2026b}, errors partition into two categories:
\emph{handled errors}, which are resolved at level $L_i$ and never
reach $L_{i+1}$, and \emph{unhandled exceptions}, which represent
genuinely irreducible failures that escalate. The error rate seen by
$L_{i+1}$ is not a compounded Bernoulli product of lower-level error
rates but the rate of genuinely irreducible errors: failures that
persist across sufficient variation to pass the threshold $\theta_i$.
The thresholds $\theta_i$ are independent by construction: each is
determined by $C_i$ and $\varepsilon$, not by the error rates of
adjacent levels.

\subsection{Deriving the Learning Rate Hierarchy}

The learning rate at level $L_i$ is inversely related to the rigor
threshold. To see why, consider the sample complexity. Each trial of
action $a$ yields a binary outcome: adaptive (success) or maladaptive
(failure). Let $p$ be the true success probability. After $n$
independent trials, the empirical success rate $\hat{p}$ satisfies
Hoeffding's inequality \cite{Hoeffding1963}:
\[
P(|\hat{p} - p| \geq \delta) \leq 2\exp(-2n\delta^2)
\]
for deviation tolerance $\delta > 0$. The entity requires confidence
$\theta_i$ that $a$ is adaptive, i.e., that $\hat{p}$ is within
$\delta$ of the true rate with probability at least $\theta_i$. Setting
$2\exp(-2n\delta^2) \leq 1 - \theta_i = \varepsilon / C_i$ and solving:
\[
n_i \;\geq\; \frac{\log(2 C_i / \varepsilon)}{2\delta^2}
\]
The effective learning rate, the fraction of experience that results
in encoding, is inversely proportional to $n_i$:
\[
\eta_i \;\propto\; \frac{1}{n_i}
\;\propto\; \frac{\delta^2}{\log(C_i / \varepsilon)}
\]
The deviation tolerance $\delta$ is level-independent: it measures how
close the empirical success rate must be to the true rate before the
entity trusts the estimate---a property of measurement accuracy, not
of the stakes. What varies across levels is $\theta_i$ (how
\emph{confident} the entity must be that it has achieved that
accuracy), which is driven by $C_i$. For fixed $\delta$, the
scaling is $\eta_i \propto 1 / \log(C_i / \varepsilon)$.

\textbf{Remark (level-independence of $\delta$).}
The assumption that $\delta$ is constant across levels is a simplifying
choice, not a necessary one. A rational entity facing lineage
extinction ($L_0$) might demand tighter measurement precision than one
facing a single failed action ($L_n$), i.e., $\delta_0 < \delta_n$.
This would make $\eta_{\mathrm{evo}}$ even smaller relative to
$\eta_{\mathrm{ind}}$, strictly strengthening the hierarchy.
Level-independent $\delta$ is therefore the conservative assumption:
it produces the \emph{weakest} version of the learning rate separation.

For $L_n$ (active context): $C_n$ is small but finite,
$\theta_n = 1 - \varepsilon / C_n$ is meaningfully less than $1$, and
$\eta_{\mathrm{ind}}$ is large. The entity updates rapidly from each
prediction-error signal.

For $L_0$ (genomic): $C_0 \gg \varepsilon$, so $\theta_0 \approx 1$
and $\eta_{\mathrm{evo}} \approx \varepsilon / C_0 \approx 0$. The
lineage updates only when certainty is overwhelming.

Since $\theta_0 \gg \theta_n$, it follows that
$\eta_{\mathrm{evo}} \ll \eta_{\mathrm{ind}}$. The learning rate
difference between evolution and individual learning is \emph{derived}
from the cost-of-failure structure, not assumed as a design parameter.

\begin{definition}[Certainty]
\label{def:certain}
Action $a$ is \emph{certain} when $U(w, K_a) \approx 0$ across
sufficient variation in $w$ to meet the threshold $\theta_i$ for
encoding level $L_i$. Formally: $P(\text{$a$ is adaptive} \mid K)
\geq \theta_i$, where the probability is taken over the distribution
of world states relevant to $L_i$.
\end{definition}

\begin{corollary}[Bounded Adaptation]
\label{cor:bounded}
An action $a$ is \emph{certain} at level $L_i$ when $U(w, K_a)
\approx 0$ across sufficient variation to meet threshold $\theta_i$.
Once $\theta_i$ is met, the expected benefit of encoding at $L_i$
exceeds the expected cost of error. The action is encoded at $L_i$ and
no longer consumes $L_n$ resources, freeing active context for the
next epistemic frontier. This encoding is durable, not literally
irreversible: environmental shift can render it maladaptive, but
correction requires effort proportional to $r_i$.
\end{corollary}

% ------------------------------------------------------------------
\section{Engagement with Prior Work}
\label{sec:prior}
% ------------------------------------------------------------------

\textbf{Dupoux, LeCun, and Malik \cite{DupouxLeCunMalik2026}.}
Their bilevel optimization framework, inner loop as development,
outer loop as evolution, captures the correct structural
relationship between encoding levels. However, they treat the learning
rate difference between levels as a design parameter: the inner loop
is fast because it is designed to be fast, and the outer loop is slow
because mutation and selection are slow. Theorem~\ref{thm:encoding}
derives this relationship from first principles: $\eta_{\mathrm{evo}}
\ll \eta_{\mathrm{ind}}$ because $C_0 \gg C_n$, not because of any
design choice.

Their framework includes a ``System~M'' that is ``not learned through
gradient descent'' but fixed by evolution. Our framework explains
\emph{why}: the rigor threshold $\theta_0 \approx 1$ means that any
component encoded at $L_0$ must be near-certain across all
environments in the species distribution. Only mechanisms that
generalize across the full environmental distribution meet this
threshold. System~M is not excluded from learning by fiat; it
is excluded because the cost of a bad learning mechanism at $L_0$ is
lineage extinction, and the resulting threshold $\theta_0$ permits
only components validated across maximal variation.

\textbf{Hutter \cite{Hutter2005}: AIXI.}
AIXI constructs the theoretically optimal agent by enumerating all
computable hypotheses, weighted by Kolmogorov complexity. The
construction assumes unbounded computation, precisely the
constraint our framework takes as central. Under bounded active
context $C_n$, the AIXI solution is unreachable: the entity cannot
enumerate hypotheses but must descend with finite resources. Our
Theorem~\ref{thm:gradient} characterizes what bounded agents actually
do: local monotone descent on $H(W \mid K)$ with convergence to local
minima, not global optimization over hypothesis space.

\textbf{Chollet \cite{Chollet2019}: Intelligence as
skill-acquisition efficiency.}
Chollet's measure of intelligence, the rate at which an agent
acquires new skills from limited experience, is structurally close
to our learning rate $\eta_M$. The key divergence is in the treatment
of prior knowledge. Chollet treats priors as a confound to be
controlled for: intelligence is efficiency \emph{net of} prior
knowledge. In our framework, prior knowledge $K$ is not a confound
but a constitutive element: $\eta_M$ operates on $K$, and the
capacity of $M$ is inseparable from the structure of $K$
\cite{McCarthy2026c}. An entity with richer $K$ does not have
``unfair advantage.'' It has done more descent.

\textbf{Schmidhuber \cite{Schmidhuber2010}: Curiosity as compression
progress.}
Schmidhuber defines intrinsic motivation as the rate of compression
progress: the reduction in description length of the agent's model
of the world. This is structurally identical to $\eta_M$: both
measure the rate of uncertainty reduction. The difference is
ontological. Schmidhuber treats curiosity as a design objective,
something to be engineered into agents. Our framework suggests that
$\gamma > 0$ (the tendency to engage $M$ even without immediate survival
pressure) is not a design choice but a necessary condition for
persistence, a claim established fully in Paper~5.

\textbf{Hinton and Nowlan \cite{Hinton1987}: Learning guiding
evolution.}
Hinton and Nowlan demonstrate via simulation that individual learning
can accelerate evolutionary search: the Baldwin effect. An
organism that can learn a trait during its lifetime effectively
smooths the fitness landscape for evolution. Our framework formalizes
this: individual learning at $L_n$ discovers actions via rapid
evaluation-driven descent ($\eta_{\mathrm{ind}}$ large). Encoding permanence
(Theorem~\ref{thm:encoding}) determines which of these actions
promote to $L_0$: only those meeting $\theta_0 \approx 1$, actions
validated across sufficient environmental variation to warrant
heritable encoding. The Baldwin effect is not an empirical curiosity
but a necessary consequence of the encoding hierarchy.

\textbf{Kirkpatrick et al.\ \cite{Kirkpatrick2017}: Catastrophic
forgetting and elastic weight consolidation.}
Kirkpatrick et al.\ demonstrate that neural networks trained
sequentially on multiple tasks catastrophically forget earlier tasks
unless protected by a consolidation mechanism. Their solution,
elastic weight consolidation (EWC), assigns higher ``importance''
to weights encoding well-established knowledge, penalizing changes to
those weights during subsequent learning. This is precisely the
encoding hierarchy: weights at lower encoding levels (higher
$\theta_i$, higher cost of error) are protected; weights at higher
levels (lower $\theta_i$) are free to change. EWC's Fisher information
matrix empirically estimates our cost-of-failure $C_i$.
The result provides direct empirical evidence that the encoding
hierarchy is not merely theoretical: artificial systems that violate
it (flat learning rate across all parameters) fail, and systems that
approximate it (EWC, progressive neural networks) succeed.

\textbf{Murray et al.\ \cite{Murray2014}: Cortical temporal
hierarchy.}
Murray et al.\ establish that intrinsic timescales of neural
activity increase systematically along the cortical hierarchy: sensory
cortex fluctuates rapidly, while prefrontal cortex maintains stable
representations over seconds. This gradient of temporal stability maps
directly onto the encoding hierarchy: fast-updating sensory areas
operate at $L_n$ (low cost of error, high reversibility), while
slow-updating prefrontal areas operate at lower encoding levels (high
cost of error, low reversibility). The empirical timescale gradient
matches the theoretical prediction $\eta_i \propto 1/\log(C_i /
\varepsilon)$: regions with higher cost of error update more slowly.

\textbf{Finn et al.\ \cite{FinnEtAl2017} and Thrun \cite{Thrun1998}:
Meta-learning.}
Model-Agnostic Meta-Learning (MAML) learns an initialization from
which task-specific learning converges rapidly. ``Learning to learn''
\cite{Thrun1998} frames this as optimizing the learning algorithm
itself. In our framework, meta-learning is $M$ operating on $M$:
the $M^{(2)}$ level. The meta-learner sits at a higher encoding level
than the base learner: it updates more slowly ($\eta_{M^{(2)}} <
\eta_M$) because the cost of a bad meta-learning rule ($C_{M^{(2)}}$)
exceeds the cost of a bad base-level update ($C_M$). The MAML inner
loop is $L_n$; the outer loop is $L_{n-1}$. The hierarchy extends:
evolution is meta-meta-$\ldots$-learning at $L_0$.

% ------------------------------------------------------------------
\section{Conclusion}
\label{sec:conclusion}
% ------------------------------------------------------------------

Monotone descent on uncertainty is not a metaphor for learning. It is
what survival-driven updating \emph{is}. Any entity with $\gamma > 0$
that persists under bounded context necessarily engages in
evaluation-driven state revision: executing actions, receiving
prediction-error signals, and updating knowledge to reduce $H(W \mid K)$.
The process is not gradient descent---there is no differentiable loss
surface over a set of propositions---but the descent property is real:
each informative update reduces expected uncertainty by exactly the
information gain. The convergence argument is scalar, with no
dependence on the dimensionality of the entity's internal
parameterization. Under positive novelty rate ($\nu(E) > 0$), descent
never terminates: novel states ensure actions with positive information
gain always exist.

The encoding hierarchy is not a design choice but a necessary
consequence of cost-of-failure differences across encoding levels.
The threshold $\theta_i = 1 - \varepsilon / C_i$ is the break-even
point: the certainty at which expected benefit of encoding exceeds
expected cost of error. From
this single equation, the entire learning rate hierarchy follows:
$\eta_{\mathrm{evo}} \ll \eta_{\mathrm{ind}}$ because
$C_0 \gg C_n$, not because evolution is ``inherently slow'' or
individual learning is ``inherently fast.''

Together with graph necessity \cite{McCarthy2026a}, convergence to
escalation \cite{McCarthy2026b}, and K/A inseparability \cite{McCarthy2026c},
this paper completes the structural account of
how knowledge is acquired, validated, and encoded. What remains is to
establish that $\gamma > 0$ is itself necessary for persistence: that
agency is not optional but required. This is the subject of Paper~5.

\section*{Acknowledgments}

The author used Claude (Anthropic) for drafting assistance, literature review, and \LaTeX{} preparation. All theoretical content, proofs, and arguments are the author's own.

\bibliographystyle{plain}
\begin{thebibliography}{99}

\bibitem{Chollet2019}
F.~Chollet.
\newblock On the measure of intelligence.
\newblock \emph{arXiv preprint arXiv:1911.01547}, 2019.

\bibitem{Cover2006}
T.~M. Cover and J.~A. Thomas.
\newblock \emph{Elements of Information Theory}.
\newblock Wiley-Interscience, 2nd edition, 2006.

\bibitem{DupouxLeCunMalik2026}
E.~Dupoux, Y.~LeCun, and J.~Malik.
\newblock Why {AI} systems don't learn and what to do about it: Lessons on
  autonomous learning from cognitive science.
\newblock \emph{arXiv preprint arXiv:2603.15381}, 2026.

\bibitem{FinnEtAl2017}
C.~Finn, P.~Abbeel, and S.~Levine.
\newblock Model-agnostic meta-learning for fast adaptation of deep networks.
\newblock In \emph{Proceedings of the 34th International Conference on Machine
  Learning (ICML)}, pages 1126--1135, 2017.

\bibitem{Hinton1987}
G.~E. Hinton and S.~J. Nowlan.
\newblock How learning can guide evolution.
\newblock \emph{Complex Systems}, 1(3):495--502, 1987.

\bibitem{Hoeffding1963}
W.~Hoeffding.
\newblock Probability inequalities for sums of bounded random variables.
\newblock \emph{Journal of the American Statistical Association},
  58(301):13--30, 1963.

\bibitem{Hutter2005}
M.~Hutter.
\newblock \emph{Universal Artificial Intelligence: Sequential Decisions Based
  on Algorithmic Probability}.
\newblock Springer, 2005.

\bibitem{Kirkpatrick2017}
J.~Kirkpatrick, R.~Pascanu, N.~Rabinowitz, J.~Veness, G.~Desjardins,
  A.~A. Rusu, K.~Milan, J.~Quan, T.~Ramalho, A.~Grabska-Barwinska,
  D.~Hassabis, C.~Summerfield, T.~Lillicrap, and D.~Kumaran.
\newblock Overcoming catastrophic forgetting in neural networks.
\newblock \emph{Proceedings of the National Academy of Sciences},
  114(13):3521--3526, 2017.

\bibitem{McCarthy2026a}
P.~D. McCarthy.
\newblock Graph structure is necessary for information preservation under
  bounded context.
\newblock Working paper, 2026.

\bibitem{McCarthy2026b}
P.~D. McCarthy.
\newblock Convergence of control architectures to escalation under bounded context.
\newblock Working paper, 2026.

\bibitem{McCarthy2026c}
P.~D. McCarthy.
\newblock Knowledge and action are inseparable under survival pressure.
\newblock Working paper, 2026.

\bibitem{Murray2014}
J.~D. Murray, A.~Bernacchia, D.~J. Freedman, R.~Romo, J.~D. Wallis,
  X.~Cai, C.~Padoa-Schioppa, T.~Pasternak, H.~Seo, D.~Lee, and
  X.-J. Wang.
\newblock A hierarchy of intrinsic timescales across primate cortex.
\newblock \emph{Nature Neuroscience}, 17(12):1661--1663, 2014.

\bibitem{Schmidhuber2010}
J.~Schmidhuber.
\newblock Formal theory of creativity, fun, and intrinsic motivation
  (1990--2010).
\newblock \emph{IEEE Transactions on Autonomous Mental Development},
  2(3):230--247, 2010.

\bibitem{Thrun1998}
S.~Thrun and L.~Pratt.
\newblock \emph{Learning to Learn}.
\newblock Kluwer Academic Publishers, 1998.

\end{thebibliography}

\end{document}
